\documentclass[11pt]{article}
\usepackage{problem_set_style}
\usepackage{kotex}
\usepackage{booktabs}
\usepackage{array}
\begin{document}
\heading{16}{샘플}{}{Spring 2027}{기말 프로젝트 샘플 (Abstract · 보고서 · 동료평가)}

\noindent{\small 아래는 \textbf{가상의 팀}이 쓴 샘플입니다. 수치는 형식을 보이기 위한 예시이며 실제 실험 결과가 아닙니다. 형식과 분량의 기준으로만 참고하세요.}

\section*{① Abstract 샘플 (기말 2주차 제출, 1쪽)}
\noindent\textbf{제목}: Fashion-MNIST에서 작은 CNN과 GBDT 비교: 딥러닝이 필요한가 \\
\textbf{팀}: 김OO, 이OO, 박OO

\vspace{4pt}
\noindent\textbf{문제 정의.} 행 하나는 \textbf{옷 사진 한 장}(28$\times$28 흑백), 라벨은 10종류 중 하나이다. 클래스가 균형이라 주 지표는 정확도와 macro-F1이다.

\noindent\textbf{왜 딥러닝인가 (초안).} 픽셀 784개를 그대로 표처럼 넣으면 이웃 픽셀의 구조를 모델이 모른다. 합성곱이 이 구조를 쓰는지가 가설이며, GBDT와의 차이로 검증한다.

\noindent\textbf{데이터와 분할.} Fashion-MNIST 70{,}000장. 공식 train 60{,}000장을 50{,}000/10{,}000으로 나눠 train/validation으로 쓰고, 공식 test 10{,}000장은 마지막에 한 번만 쓴다.

\noindent\textbf{베이스라인.} 픽셀 784개를 입력으로 한 \textbf{GBDT}. 같은 train/validation으로 튜닝한다.

\noindent\textbf{계획.} 합성곱 2층의 작은 CNN을 만들고, 에포크 곡선으로 조기 종료 지점을 정한다. 데이터 증강(좌우 반전)과 드롭아웃의 효과를 구성 요소 제거 실험으로 확인한다.

\newpage
\section*{② 보고서 샘플 (기말 6주차 제출, 축약본)}
\noindent\textbf{1. 문제 정의.} 사진 한 장이 행, 라벨은 10종류. 주 지표는 정확도와 macro-F1. 베이스라인은 GBDT(픽셀 입력).

\noindent\textbf{2. 데이터와 전처리.} 공식 test는 한 번만 사용했다. 픽셀 값은 0--1로 스케일링(데이터 전체 상수 255로 나눔, 통계량 누수 없음).

\noindent\textbf{3. 모델 선택과 이유.} 합성곱이 이웃 픽셀의 모양을 쓰는지 보기 위해 CNN을, 구조를 모르는 비교 대상으로 GBDT를 골랐다.

\begin{center}\small
\begin{tabular}{@{}lcc@{}}\toprule
모델 (validation) & 정확도 & 비고 \\\midrule
GBDT (픽셀 784) & 0.89 & 베이스라인 \\
CNN 2층 & 0.91 & \\
CNN 2층 + 데이터 증강 & 0.92 & \\
CNN 2층 + 증강 + 드롭아웃 & 0.93 & 에포크 15에서 종료 \\\bottomrule
\end{tabular}
\end{center}

\noindent\textbf{4. 결과.} 최종 test는 한 번만: CNN 0.92, GBDT 0.88. 학습 곡선(에포크 단위)에서 에포크 15 이후 validation이 내려가기 시작해 15에서 멈췄다.

\noindent\textbf{5. 수식적 정당화와 왜 딥러닝인가.} 합성곱은 파라미터를 위치마다 공유해 3$\times$3 필터 하나로 어디서든 같은 모양을 찾는다. 같은 크기의 완전연결층보다 파라미터가 적다. 그 결과 GBDT보다 test 정확도가 4\%p 높았다. 다만 정확도 차이가 크지 않아, 표 데이터였다면 GBDT로 충분했을 것이다.

\noindent\textbf{6. 재현 방법.} 공식 Fashion-MNIST, 시드 0, 배치 16, 에포크 15, 학습률 0.001. \texttt{run.sh} 한 줄로 다시 실행된다.

\noindent\textbf{7. 한계와 잘 안 된 부분.} 배치 정규화를 넣었더니 validation이 오르락내리락해서 뺐다. 배치 크기가 작아(16) 통계가 불안정했다고 진단했다.

\newpage
\section*{③ 동료평가 샘플 (기말 8주차 제출)}
\noindent\textbf{평가 대상}: 다른 팀의 ``영화 리뷰 감성 분류'' 보고서 \quad \textbf{작성자}: 이OO

\begin{center}\small
\begin{tabular}{@{}p{3.4cm}p{9cm}@{}}\toprule
항목 & 평가 \\\midrule
주장의 명확성 & LSTM이 TF-IDF보다 낫다는 주장이 분명하다. \\
비교 대상 & TF-IDF + 로지스틱회귀 베이스라인이 없다. 왜 딥러닝인가가 정확도가 높아서 한 줄이다. \\
확장성 & 영화 리뷰만 다뤘다. 다른 도메인(상품 리뷰)에서의 확인이 없다. \\
재현성 & 데이터 출처, 시드, 에포크 수, 실행 명령이 모두 있어 다시 해 볼 수 있다. 좋다. \\
한계 & 긴 문장에서 성능이 떨어지는 것을 보이고 원인을 진단했다. 좋다. \\
\bottomrule\end{tabular}\end{center}

\noindent\textbf{반박 가능한 질문.} TF-IDF + 로지스틱회귀 같은 단순한 베이스라인을 돌렸나요? 그 정확도와 LSTM의 차이가 계산 비용을 감당할 만큼인가요?

\end{document}
